\documentclass{article}

\usepackage[T1]{fontenc}
\usepackage{lmodern}
\usepackage{amsmath}
\usepackage{amssymb}
\usepackage{amsthm}
\usepackage{mathtools}
\usepackage[hidelinks]{hyperref}

\newtheorem*{mainproposition}{Proposition}
\newtheorem{lemma}{Lemma}

\DeclareMathOperator{\wt}{wt}

\title{A polynomial run-exit bound}
\author{}
\date{Self-contained hand proof, 11 November 2026}

\begin{document}

\maketitle

\section{Main proposition}\label{sec:main}

\begin{mainproposition}[weight at least seventeen; independent review\slash formalization pending]
  Let \((a_i,b_i)_{i=0}^{18}\) be polynomial pairs over \(\mathbb{F}_2\), each nonzero polynomial
  having degree at most 18. Assume either of these exact polynomial systems:

  \textbf{Forward system F}, in \(\mathbb{F}_2[y]\):
  \begin{equation}\label{eq:forward}
    a_i = (y + y^2) a_{i-2} + b_{i-1} + b_{i-2} + y^2 b_i \quad (2 \le i \le 18).
  \end{equation}

  \textbf{Backward system B}, in \(\mathbb{F}_2[x]\):
  \begin{equation}\label{eq:backward}
    (x + x^2) b_i = a_i + x^2 a_{i-1} + x^2 a_{i-2} + b_{i-2} \quad (2 \le i \le 18).
  \end{equation}

  Let \(\wt(p)\) count the nonzero coefficients of \(p\), with \(\wt(0) = 0\), and put
  \begin{equation}
    w_i = \wt(a_i) + \wt(b_i), \quad W = \sum_{i=0}^{18} w_i.
  \end{equation}
  Assume adjacent nonemptiness:
  \begin{equation}\label{eq:adjacent}
    w_i + w_{i+1} \ge 1 \quad (0 \le i \le 17).
  \end{equation}
  Let \(z\) count indices with \(w_i = 0\), and set \(E = W - (19 - z)\). For F define
  \(\varepsilon = 1\) exactly when \(w_{16} = w_{18} = 0\); for B define \(\varepsilon = 1\) exactly
  when \(w_0 = w_2 = 0\). Otherwise set \(\varepsilon = 0\). Then
  \begin{equation}\label{eq:main-bound}
    E \ge z - 1 - \varepsilon, \quad W \ge 18 - \varepsilon \ge 17.
  \end{equation}
\end{mainproposition}

\textbf{Status and scope.} All results in this draft are PROPOSITION with the named gap
independent review\slash formalization pending, not refereed THEOREM. This is a hand proof from the
displayed hypotheses alone, with no charge, code-model, matrix-recognition, or stabilizer
dependency. Both recurrences receive separate algebraic proofs. No enumeration, LP, search
certificate, or computer proof is used. There is no polynomial quotient, periodic identification,
or truncation of products. The degree cap is not used in the argument, so the same result holds
for arbitrary finite polynomial slices satisfying these equations and the same index range.

\textbf{Proof organization.} Lemmas~\ref{lem:branch}--\ref{lem:four} establish directional
branching and local gap costs. Lemmas~\ref{lem:propagation}--\ref{lem:exit-alternative} establish
short-run propagation and the run-exit alternative. Section~\ref{sec:accounting} partitions the
excess without double counting and proves the main proposition. Every needed local lemma is
proved here. This document makes no claim of a bound twenty.

\section{Polynomial facts and a branch between zeros}\label{sec:polynomial}

For the local calculations only, rename the indeterminate in either system as \(T\), and put
\(t = T + T^2\). This renames a variable, not a recurrence. A slice is empty if \(w_i = 0\), a
singleton if \(w_i = 1\), and heavy if \(w_i \ge 2\). A singleton pair has exactly one
nonzero block, which is a monomial. In particular \(u + T^2 v\) and \(T^2 u + v\) each have weight
one whenever \((u,v)\) is a singleton.

If \(\rho(p)\) counts maximal consecutive runs of exponents in the support of \(p\), with
\(\rho(0) = 0\), then
\begin{equation}\label{eq:runs}
  \wt(tp) = \wt((1 + T)p) = 2\rho(p).
\end{equation}
Indeed a run from exponent \(c\) to \(d\) contributes exactly \(T^c + T^{d+1}\) to \((1 + T)p\).
Endpoints from distinct runs are distinct. A monomial shift does not change weight. Thus
\(\wt(tp) \ge 2\) for \(p \ne 0\). We also use \(\wt(u + v) \le \wt(u) + \wt(v)\) and the fact
that \(p(1)\) is the parity of \(\wt(p)\). Both follow coefficient by coefficient over
\(\mathbb{F}_2\). The polynomial ring is a domain, so nonzero factors can be cancelled.

\begin{lemma}[directional branch]\label{lem:branch}
  Suppose \(1 \le r \le 17\) and slices \(r - 1, r + 1\) are empty. In system F, \(b_r = 0\); if
  \(r \le 16\) and slice \(r\) is nonempty, then \(w_{r+2} \ge 2\). In system B, \(a_r = 0\); if
  \(r \ge 2\) and slice \(r\) is nonempty, then \(w_{r-2} \ge 2\).
\end{lemma}

\begin{proof}[Forward proof]
  Applying equation~\eqref{eq:forward} at \(r + 1\) gives \(0 = b_r\). At \(r + 2\), when this
  index is at most 18, it gives
  \begin{equation}
    a_{r+2} + T^2 b_{r+2} = t a_r.
  \end{equation}
  The sum of the two physical block weights is at least the weight of their sum, hence at least
  \(2\rho(a_r) \ge 2\) when \(a_r \ne 0\).
\end{proof}

\begin{proof}[Backward proof]
  Applying equation~\eqref{eq:backward} at \(r + 1\) gives \(0 = T^2 a_r\), so \(a_r = 0\).
  Applying equation~\eqref{eq:backward} at \(r\), when \(r \ge 2\), gives
  \begin{equation}
    T^2 a_{r-2} + b_{r-2} = t b_r.
  \end{equation}
  The same coefficient-weight inequality gives \(w_{r-2} \ge 2\rho(b_r) \ge 2\). This proves both
  assertions.
\end{proof}

There is exactly one possible exceptional middle index per direction: \(r = 17\) in F would
require equation 19, and \(r = 1\) in B would require equation 1. Neither equation is among the
hypotheses. No branch cost is asserted at these exceptions.

\section{Local costs inside gaps}\label{sec:gaps}

\begin{lemma}[two occupied slices between zeros]\label{lem:two}
  If \(0 \le k \le 15\), slices \(k, k + 3\) are empty, and the intermediate slices are not both
  empty, then
  \begin{equation}
    w_{k+1} + w_{k+2} \ge 4.
  \end{equation}
\end{lemma}

\begin{proof}[Forward proof]
  Write the intermediate slices as \((p,q), (u,v)\). Equations at \(k + 2, k + 3\) are
  \begin{equation}
    u = q + T^2 v, \quad tp + q + v = 0.
  \end{equation}
  Evaluation at \(T = 1\) gives \(u(1) = 0\). If \(u \ne 0\), then \(\wt(u) \ge 2\). Also
  \(\wt(p) + \wt(q) + \wt(v) \ge 2\): weight zero forces \(u = 0\), and weight one leaves a single
  nonzero term in \(tp + q + v = 0\). Thus the total is at least four. If \(u = 0\), then
  \(q = T^2 v\) and \(tp = (1 + T^2)v\). Cancelling \(1 + T\) yields \(Tp = (1 + T)v\). Evaluation
  at zero shows that \(v = Th\) for a polynomial \(h\); then \(p = (1 + T)h\) and \(q = T^3 h\).
  Nonemptiness gives \(h \ne 0\), and total weight is \(2\wt(h) + 2\rho(h) \ge 4\).
\end{proof}

\begin{proof}[Backward proof]
  The equations for the same named slices are instead
  \begin{equation}
    tv = u + T^2 p, \quad q = T^2 (u + p).
  \end{equation}
  Evaluation at one gives \(q(1) = 0\). If \(q \ne 0\), its weight is at least two. The remaining
  weight \(\wt(p) + \wt(u) + \wt(v)\) is at least two, since weight zero forces \(q = 0\) and
  weight one cannot satisfy \(tv + u + T^2 p = 0\). If \(q = 0\), then \(u = p\) and
  \(tv = (1 + T^2)p\). Cancelling \(1 + T\) gives \(Tv = (1 + T)p\). Hence \(p = Th\), \(u = Th\),
  \(v = (1 + T)h\) for \(h \ne 0\). The total is again \(2\wt(h) + 2\rho(h) \ge 4\).
\end{proof}

\begin{lemma}[three occupied slices between zeros]\label{lem:three}
  If \(0 \le k \le 14\), slices \(k, k + 4\) are empty, and all three intermediate slices are
  occupied, then
  \begin{equation}
    w_{k+1} + w_{k+2} + w_{k+3} \ge 4.
  \end{equation}
\end{lemma}

\begin{proof}[Forward proof]
  A failure would make all three pairs \((p,q), (u,v), (h,l)\) singletons. The equations at
  \(k + 2, k + 3, k + 4\) are
  \begin{equation}
    u = q + T^2 v, \quad h = tp + v + q + T^2 l, \quad 0 = tu + l + v.
  \end{equation}
  If \(q = 0\), the first equation forces both \(u, v\) to vanish or both to be nonzero,
  contradicting the second singleton. Thus \(q\) is a monomial and \(p = 0\). If \(u \ne 0\), then
  \(v = 0\) and the last equation gives \(l = tu\) of weight two, contradicting the third
  singleton. Consequently \(u = 0\), \(v\) is a monomial, and \(l = v\), so \(h = 0\). The middle
  equation now gives \(q = (1 + T^2)v\), of weight two, a contradiction.
\end{proof}

\begin{proof}[Backward proof]
  The equations for the three singleton slices are
  \begin{equation}
    tv = u + T^2 p, \quad tl = h + T^2 u + T^2 p + q, \quad v = T^2 (h + u).
  \end{equation}
  If \(v \ne 0\), it is a monomial and \(u = 0\), so the first equation forces \(T^2 p = tv\) of
  weight two, impossible in the first singleton. Hence \(v = 0\) and \(u\) is a monomial. The last
  equation gives \(h = u\), so \(l = 0\). The first gives \(T^2 p = u\), so \(q = 0\). The middle
  equation is then \(0 = u + T^2 u + u = T^2 u\), a contradiction.
\end{proof}

\begin{lemma}[four occupied slices next to a zero]\label{lem:four}
  In system F, an empty slice \(k\), with \(0 \le k \le 14\), followed by four occupied slices
  satisfies
  \begin{equation}
    \sum_{i=k+1}^{k+4} w_i \ge 5.
  \end{equation}
  In system B, an empty slice \(R\), with \(4 \le R \le 18\), preceded by four occupied slices
  satisfies
  \begin{equation}
    \sum_{i=R-4}^{R-1} w_i \ge 5.
  \end{equation}
  Neither assertion requires another empty slice at the other end.
\end{lemma}

\begin{proof}[Forward proof]
  If all four slices \((p,q), (u,v), (h,l), (m,n)\) are singletons in increasing index order, the
  equations at \(k + 2, k + 3, k + 4\) are
  \begin{equation}
    u + T^2 v = q, \quad h + T^2 l = tp + v + q, \quad m + T^2 n = tu + l + v.
  \end{equation}
  The first left side is a monomial, so \(q\) is a monomial and \(p = 0\). If the second slice
  were B-only, then \(u = 0, q = T^2 v\), and the second equation would give
  \(h + T^2 l = (1 + T^2)v\) of weight two, a contradiction. Thus the second slice is A-only, with
  \(u = q, v = 0\). The third equation to solve for a slice is now \(h + T^2 l = q\). If the third
  slice is A-only, then \(h = q, l = 0\), and the fourth-slice equation has right side \(tq\), of
  weight two. If it is B-only, then \(h = 0, q = T^2 l\), and that right side is
  \(tq + l = (1 + T^3 + T^4)l\), of weight three. Both contradict the fourth singleton, whose left
  side \(m + T^2 n\) is a monomial.
\end{proof}

\begin{proof}[Backward proof]
  Name singleton slices \((p,q), (u,v), (h,l), (m,n)\) in decreasing order at
  \(R - 1, R - 2, R - 3, R - 4\). The equations at \(R, R - 1, R - 2\) are
  \begin{equation}
    T^2 u + v = T^2 p, \quad T^2 h + l = tq + p + T^2 u, \quad T^2 m + n = tv + u + T^2 h.
  \end{equation}
  The first left side is a monomial, so \(p\) is a monomial and \(q = 0\). If the second slice
  were A-only, then \(u = p, v = 0\), and the second equation would give
  \(T^2 h + l = (1 + T^2)p\) of weight two. Thus the second slice is B-only, with
  \(u = 0, v = T^2 p\). The next equation is \(T^2 h + l = p\). If the third slice is B-only, then
  \(h = 0, l = p\), and the last right side is \(t T^2 p\), of weight two. If it is A-only, then
  \(l = 0, T^2 h = p\), and the last right side is \((1 + T^3 + T^4)p\), of weight three. Neither
  is the monomial \(T^2 m + n\) required by the fourth singleton.
\end{proof}

In a gap of length five between successive zero slices, all four interior slices are occupied.
Lemma~\ref{lem:four} applies from the left endpoint in F and from the right endpoint in B. These
endpoints satisfy \(k \le 13\) and \(R \ge 5\), respectively. Thus such a gap has interior weight
at least five. All monomials counted in the proofs have distinct exponents; none may be discarded
if an exponent exceeds 18. Such an identity would violate the hypotheses rather than supply a
further cancellation.

\section{Short runs and exit classification}\label{sec:runs}

Orient the list of zero slices in the forcing direction: increasing index for F and decreasing
index for B. A \emph{short gap} is a gap of length two between successive zero slices. A
\emph{short run} is a maximal sequence of \(m \ge 1\) consecutive short gaps. Its \(m\) middle
slices are nonempty by \eqref{eq:adjacent}. The \emph{exit target} is the next slice after its
last zero in the forcing direction. It either lies inside the following longer gap, lies in the
tail after the last zero, or falls outside the index range.

\begin{lemma}[exact propagation and exceptional exits]\label{lem:propagation}
  Write the active polynomials in the \(m\) middle slices of a short run as
  \(P_1, \dots, P_m\) in forcing order. In F these are A-only slices and in B they are B-only
  slices. Renaming the indeterminate as \(T\) and writing \(t = T + T^2\), one has
  \begin{equation}\label{eq:run-powers}
    P_j = t^{j-1} P_1 \quad (1 \le j \le m), \quad P_1 \ne 0.
  \end{equation}
  At an on-board exit \(e\), respectively,
  \begin{equation}\label{eq:exit-feed}
    a_e + T^2 b_e = t P_m \quad \text{(F)}, \quad
    T^2 a_e + b_e = t P_m \quad \text{(B)}.
  \end{equation}
  Thus every middle after the first, and every on-board exit, has weight at least two. There is
  at most one missing exit: the final middle 17 in F or the final middle 1 in B. Its occurrence is
  exactly \(\varepsilon = 1\) in the main statement.
\end{lemma}

\begin{proof}
  Lemma~\ref{lem:branch} gives the active block of every middle. In F, at the next middle
  \(r + 2\), the B terms at \(r, r + 1, r + 2\) vanish; hence its A polynomial is \(t\) times that
  at \(r\). At an exit the same equation gives the F identity in \eqref{eq:exit-feed}. In B, the
  equation at a middle \(r\) is \(t b_r = T^2 a_{r-2} + b_{r-2}\) because \(a_r = a_{r-1} = 0\). If
  \(r - 2\) is the next middle, its A polynomial also vanishes, proving the B propagation step;
  otherwise it gives the B exit identity. Induction proves \eqref{eq:run-powers}. The domain
  property and \eqref{eq:runs} show weight at least two at every positive power and every on-board
  exit.

  A run ending at zero 18 in F has last middle 17 and would exit at 19. Every other F run has an
  on-board exit: in particular zero 17 has target 18. In B a run ending at zero 0 has last middle 1
  and would exit at \(-1\); zero 1 still has target 0. Only the final run can lack an exit. Such a
  run exists precisely when the final zero gap is \((16,18)\) in F or \((2,0)\) in B, as specified
  by \(\varepsilon\).
\end{proof}

\begin{lemma}[cheap exit forces a feed of weight three or four]\label{lem:cheap-exit}
  Suppose three successive exit slices in forcing order have weights \(2, 1, 1\). If their feeding
  middle polynomial is \(P\), then \(\wt(P)\) is three or four. More precisely \(P\) is a monomial
  shift of \(1 + T + T^2\) or \(1 + T + T^2 + T^3\). The following statements specify the exact
  index ranges and equations.
\end{lemma}

\begin{proof}[Forward proof]
  Let the feeding middle be at \(r\), with \(1 \le r \le 14\), and zeros at \(r - 1, r + 1\). It
  is \((P, 0)\). Write \((A_i, B_i)\) for the slices at \(r + 2 + i\), for \(i = 0, 1, 2\). The
  equations at \(r + 2, r + 3, r + 4\) are
  \begin{equation}\label{eq:cheap-forward}
    \begin{gathered}
      A_0 + T^2 B_0 = tP, \quad A_1 + T^2 B_1 = B_0, \\
      A_2 + T^2 B_2 = t A_0 + B_0 + B_1.
    \end{gathered}
  \end{equation}
  The second equation makes \(B_0\) a monomial; the first slice's weight two then makes \(A_0\) a
  monomial. If the second slice is B-only, the last right side has even coefficient parity:
  \(t A_0\) has even weight and \(B_0, B_1\) are monomials. It cannot be the monomial on the left.
  Thus the second slice is A-only, with \(A_1 = B_0, B_1 = 0\). For \(t A_0 + B_0\) to be a
  monomial, \(B_0\) must cancel exactly one term of \(t A_0\). Hence \(B_0 = T A_0\) or
  \(B_0 = T^2 A_0\). Write \(A_0 = T^a\). The feed equation gives
  \begin{equation}
    P = T^{a-1}(1 + T + T^2) \quad \text{or} \quad
    P = T^{a-1}(1 + T + T^2 + T^3), \quad a \ge 1.
  \end{equation}
  To justify the division, \(a = 0\) would give a nonzero constant on the left of the feed, while
  \(T\) divides its right side. For \(a \ge 1\), cancel \(T(1 + T)\) using
  \(1 + T^3 = (1 + T)(1 + T + T^2)\) and \(1 + T^4 = (1 + T)(1 + T + T^2 + T^3)\). Both displayed
  feed forms have the asserted weight.
\end{proof}

\begin{proof}[Backward proof]
  Let the feeding middle be at \(r\), with \(4 \le r \le 17\) and zeros at \(r - 1, r + 1\). It is
  \((0, P)\). Write \((A_i, B_i)\) for the slices at \(r - 2 - i\). The equations at
  \(r, r - 1, r - 2\) give
  \begin{equation}\label{eq:cheap-backward}
    \begin{gathered}
      T^2 A_0 + B_0 = tP, \quad T^2 A_1 + B_1 = T^2 A_0, \\
      T^2 A_2 + B_2 = t B_0 + A_0 + T^2 A_1.
    \end{gathered}
  \end{equation}
  Now the second equation makes \(A_0\) a monomial and weight two of the first slice makes \(B_0\)
  a monomial. If the second slice is A-only, then \(A_1 = A_0\) and the last right side is
  \(t B_0 + (1 + T^2)A_0\), of even coefficient parity. It cannot be a monomial. Thus the second
  slice is B-only, and \(t B_0 + A_0\) must be a monomial. Necessarily \(A_0 = T B_0\) or
  \(A_0 = T^2 B_0\). Writing \(B_0 = T^b\), the feed yields
  \begin{equation}
    P = T^{b-1}(1 + T + T^2) \quad \text{or} \quad
    P = T^{b-1}(1 + T + T^2 + T^3), \quad b \ge 1.
  \end{equation}
  As above \(b = 0\) is excluded by the constant coefficient, and cancellation of \(T(1 + T)\)
  proves the formulas for \(b \ge 1\). This proves the lemma from the separate backward
  equations.
\end{proof}

\begin{lemma}[exit alternative]\label{lem:exit-alternative}
  Suppose a short run feeds a longer zero gap \(D\) of length \(d \ge 4\). Let \(P\) be its final
  middle polynomial and let \(E_D\) be the excess inside \(D\), namely the sum of \(w_i - 1\) over
  its interior. Then
  \begin{equation}\label{eq:exit-alternative}
    E_D \ge 2 \quad \text{or} \quad \wt(P) \ge 3.
  \end{equation}
  For length three, \(E_D \ge 2\) unconditionally.
\end{lemma}

\begin{proof}
  The exit target is the first interior slice and has weight at least two by
  Lemma~\ref{lem:propagation}. If \(E_D \le 1\), that slice has weight exactly two and every other
  interior slice is a singleton. There are at least three interior slices, so
  Lemma~\ref{lem:cheap-exit} applies. Its equations are tested: a gap of length at least four
  requires \(r + 5 \le 18\) in F or \(r - 5 \ge 0\) in B, stronger than the respective index
  ranges in Lemma~\ref{lem:cheap-exit}. Thus \(\wt(P) \ge 3\). For length three use
  Lemma~\ref{lem:two}.
\end{proof}

\section{Disjoint excess accounting and the main proof}\label{sec:accounting}

If \(z = 0\), all nineteen slices are occupied, giving \(W \ge 19\) and
\(E \ge 0 \ge z - 1 - \varepsilon\). Hence assume \(z \ge 1\). There are \(z - 1\)
successive-zero gaps. All have length at least two by \eqref{eq:adjacent}. Let \(G\) count short
gaps, \(M = z - 1 - G\) count longer gaps, and \(R_2\) count maximal short runs. Every interior
slice of a successive-zero gap is occupied, so its contribution to \(E\) is exactly its weight
minus one.

\textbf{Baseline inside runs.} In a short run of \(m\) gaps, reserve one excess unit at every
middle after the first, totalling \(m - 1\). These units exist by Lemma~\ref{lem:propagation}.
Across all short runs this reserves \(G - R_2\) units.

\textbf{Baseline inside longer gaps.} Reserve one excess unit in every longer gap, totalling
\(M\). Length three has at least two units by Lemma~\ref{lem:two}. Length four has at least one
by Lemma~\ref{lem:three}. Length at least five has at least one in the first four interior slices
in forcing order by Lemma~\ref{lem:four}. These longer-gap interiors are disjoint from each other
and from all short-run middles. Thus the baseline reserves \(G - R_2 + M\) units.

\textbf{An additional unit for each run that has an exit.} If a run is followed by a gap of length
three, that gap has a second excess unit after its reserved baseline. If its following gap has
length at least four, Lemma~\ref{lem:exit-alternative} says either the same second gap unit
exists, or the run's terminal middle has weight at least three. In the latter case that middle
has at least two excess units. At most one was reserved there in the run baseline: one for
\(m \ge 2\), zero for \(m = 1\). At least one further unit therefore remains. Knowing only that
the terminal middle was heavy would not suffice; Lemma~\ref{lem:cheap-exit} supplies the needed
weight at least three.

Distinct short runs have distinct terminal middles and distinct following longer gaps. Run
interiors are disjoint from longer-gap interiors. Thus none of these selected additional units is
a reserved baseline unit or a unit already selected for another run.

If instead the final run exits into an on-board tail, Lemma~\ref{lem:propagation} makes its
target heavy. Its excess unit lies outside all baseline regions and supplies that run's
additional unit. Only a run with a missing exit receives no required additional unit.
Lemma~\ref{lem:propagation} identifies precisely \(\varepsilon\) such runs. Therefore the number
of additional units is at least \(R_2 - \varepsilon\), and
\begin{equation}
  \begin{aligned}
    E &\ge (G - R_2 + M) + (R_2 - \varepsilon) \\
      &= G + M - \varepsilon = z - 1 - \varepsilon.
  \end{aligned}
\end{equation}
Finally,
\begin{equation}
  W = 19 - z + E \ge 18 - \varepsilon \ge 17.
\end{equation}
This includes \(z = 1\), where \(G = M = R_2 = \varepsilon = 0\) and \(W \ge 18\). The main
proposition is proved. \qed

\textbf{Endpoint distinction.} The only possible discount is the nonexistent target at index 19
in F or \(-1\) in B. An existing target at index 18 in F or index 0 in B still costs an excess
unit, even when no later zero exists. The result gives eighteen whenever that missing-exit
configuration is absent.

\textbf{Review boundary.} This draft contains the local seventeen result and all its ingredients.
It does not use or modify the earlier local-fifteen draft. No connection to an external model is
asserted. All claims remain PROPOSITION pending independent review and formalization.

\end{document}
